% GENERATED FILE. Edit canonical lessons, then run npm run generate:books.
% canonical-source-hash: fnv1a64:55128eb73d568ffe
% canonical-lessons: FR-C485-guerir, FR-C485-guerison, FR-C485-retablir, FR-C485-convalescence, FR-C485-soigner

\chapter{Tomber malade et guérir}
\label{ch:fr-tomber-malade-et-guerir}
\hlchaptermodality{485}

\begin{chapteropening}
\textbf{By the end of this chapter:} \emph{I can talk about getting better after an illness.}

Complete the last lesson of chapter 485: I can talk about getting better after an illness.
\end{chapteropening}

\section[guérir]{guérir — to heal, to cure; to get better}

\label{lesson:FR-C485-guerir}



\begin{quote}
Before the new one: say the French for injection, then the French for waiting room.
\end{quote}

\subsection*{You'll want to know: guérir}
\textbf{guérir} — \textquotedblleft{}to heal, to cure; to get better\textquotedblright{}.

\begin{quote}
Ma grand-mère a guéri de sa grippe en une semaine.
\end{quote}

\begin{quote}
\emph{My grandmother got over her flu in a week.}
\end{quote}

\emph{Guérir de} — to recover from. A regular \emph{-ir} verb: \emph{nous guérissons}.

\begin{cousinweb}
\emph{guérir} comes from Frankish \emph{warjan} \textquotedblleft{}to protect, to defend\textquotedblright{} (EtymDB), a Germanic word; the sense moved from \textquotedblleft{}protect\textquotedblright{} to \textquotedblleft{}heal\textquotedblright{}.
\end{cousinweb}

\begin{grammarlens}[title={What you've built}]
One more word for telling, explaining and describing at B1.
\end{grammarlens}

\subsection*{Your turn}
\begin{itemize}
\raggedright
  \item \emph{Say it:} \emph{guérir}
  \item \emph{Say it:} \emph{guérir}, once more
  \item \emph{Say it:} a sentence of your own with \emph{guérir}
  \item \emph{Recall:} say the French for injection, then the French for waiting room, then say \emph{guérir} again
\end{itemize}

\begin{tcolorbox}[breakable,colback=teal!4,colframe=teal!35!black,title={Before you move on}]
What does guérir mean? (\textquotedblleft{}To heal, to cure; to get better\textquotedblright{}.) Say it once more.
\end{tcolorbox}
\section[la guérison]{la guérison — recovery, cure}

\label{lesson:FR-C485-guerison}



\begin{quote}
Before the new one: say the French for waiting room, then the French for to heal.
\end{quote}

\subsection*{You'll want to know: la guérison}
\textbf{la guérison} — \textquotedblleft{}recovery, cure\textquotedblright{}.

\begin{quote}
Les médecins sont contents de sa guérison.
\end{quote}

\begin{quote}
\emph{The doctors are pleased with his recovery.}
\end{quote}

Built on \emph{guérir}.

\begin{cousinweb}
\emph{la guérison} is built on \emph{guérir} \textquotedblleft{}to heal\textquotedblright{}, which comes from Frankish \emph{warjan} \textquotedblleft{}to protect\textquotedblright{} (EtymDB).
\end{cousinweb}

\begin{grammarlens}[title={What you've built}]
One more word for telling, explaining and describing at B1.
\end{grammarlens}

\subsection*{Your turn}
\begin{itemize}
\raggedright
  \item \emph{Say it:} \emph{la guérison}
  \item \emph{Say it:} \emph{la guérison}, once more
  \item \emph{Say it:} a sentence of your own with \emph{la guérison}
  \item \emph{Recall:} say the French for waiting room, then the French for to heal, then say \emph{la guérison} again
\end{itemize}

\begin{tcolorbox}[breakable,colback=teal!4,colframe=teal!35!black,title={Before you move on}]
What does la guérison mean? (\textquotedblleft{}Recovery, cure\textquotedblright{}.) Say it once more.
\end{tcolorbox}
\section[rétablir]{rétablir — to restore; (se rétablir) to recover}

\label{lesson:FR-C485-retablir}



\begin{quote}
Before the new one: say the French for to heal, then the French for recovery.
\end{quote}

\subsection*{You'll want to know: rétablir}
\textbf{rétablir} — \textquotedblleft{}to restore; (se rétablir) to recover\textquotedblright{}.

\begin{quote}
Après l'opération, elle s'est vite rétablie.
\end{quote}

\begin{quote}
\emph{After the operation, she recovered quickly.}
\end{quote}

Mostly reflexive for health: \emph{se rétablir}. \emph{Rétablir le courant} — to restore the power.

\begin{cousinweb}
\emph{rétablir} is built from \emph{re-} and \emph{établir} \textquotedblleft{}to set up\textquotedblright{}, inherited from Latin \emph{stabilīre} \textquotedblleft{}to make firm\textquotedblright{}, from \emph{stabilis}, built on \emph{stāre} \textquotedblleft{}to stand\textquotedblright{} (Lewis \& Short).
\end{cousinweb}

\begin{grammarlens}[title={What you've built}]
One more word for telling, explaining and describing at B1.
\end{grammarlens}

\subsection*{Your turn}
\begin{itemize}
\raggedright
  \item \emph{Say it:} \emph{rétablir}
  \item \emph{Say it:} \emph{rétablir}, once more
  \item \emph{Say it:} a sentence of your own with \emph{rétablir}
  \item \emph{Recall:} say the French for to heal, then the French for recovery, then say \emph{rétablir} again
\end{itemize}

\begin{tcolorbox}[breakable,colback=teal!4,colframe=teal!35!black,title={Before you move on}]
What does rétablir mean? (\textquotedblleft{}To restore; (se rétablir) to recover\textquotedblright{}.) Say it once more.
\end{tcolorbox}
\section[la convalescence]{la convalescence — convalescence, recovery period}

\label{lesson:FR-C485-convalescence}



\begin{quote}
Before the new one: say the French for recovery, then the French for to restore.
\end{quote}

\subsection*{You'll want to know: la convalescence}
\textbf{la convalescence} — \textquotedblleft{}convalescence, recovery period\textquotedblright{}.

\begin{quote}
Pendant sa convalescence, il a lu beaucoup de livres.
\end{quote}

\begin{quote}
\emph{During his recovery he read a lot of books.}
\end{quote}

\emph{Être en convalescence} — to be convalescing.

\begin{cousinweb}
\emph{la convalescence} is a learned borrowing from Late Latin \emph{convalēscentia} \textquotedblleft{}regaining of health\textquotedblright{} (EtymDB), from \emph{convalēscere} \textquotedblleft{}to grow strong\textquotedblright{}, built on \emph{valēre} \textquotedblleft{}to be strong\textquotedblright{} (Lewis \& Short).
\end{cousinweb}

\begin{grammarlens}[title={What you've built}]
One more word for telling, explaining and describing at B1.
\end{grammarlens}

\subsection*{Your turn}
\begin{itemize}
\raggedright
  \item \emph{Say it:} \emph{la convalescence}
  \item \emph{Say it:} \emph{la convalescence}, once more
  \item \emph{Say it:} a sentence of your own with \emph{la convalescence}
  \item \emph{Recall:} say the French for recovery, then the French for to restore, then say \emph{la convalescence} again
\end{itemize}

\begin{tcolorbox}[breakable,colback=teal!4,colframe=teal!35!black,title={Before you move on}]
What does la convalescence mean? (\textquotedblleft{}Convalescence, recovery period\textquotedblright{}.) Say it once more.
\end{tcolorbox}
\section[soigner]{soigner — to look after, to treat, to nurse}

\label{lesson:FR-C485-soigner}



\begin{quote}
Before the new one: say the French for to restore, then the French for convalescence.
\end{quote}

\subsection*{You'll want to know: soigner}
\textbf{soigner} — \textquotedblleft{}to look after, to treat, to nurse\textquotedblright{}.

\begin{quote}
Le vétérinaire a bien soigné notre chat.
\end{quote}

\begin{quote}
\emph{The vet treated our cat well.}
\end{quote}

\emph{Se soigner} — to look after oneself. \emph{Soigné} — neat, well-groomed.

\begin{cousinweb}
\emph{soigner} is built on \emph{soin} \textquotedblleft{}care\textquotedblright{}, usually traced to a Frankish word \emph{sunnija} \textquotedblleft{}care, concern\textquotedblright{} (EtymDB).
\end{cousinweb}

\begin{grammarlens}[title={What you've built}]
One more word for telling, explaining and describing at B1.
\end{grammarlens}

\subsection*{Your turn}
\begin{itemize}
\raggedright
  \item \emph{Say it:} \emph{soigner}
  \item \emph{Say it:} \emph{soigner}, once more
  \item \emph{Say it:} a sentence of your own with \emph{soigner}
  \item \emph{Recall:} say the French for to restore, then the French for convalescence, then say \emph{soigner} again
\end{itemize}

\begin{tcolorbox}[breakable,colback=teal!4,colframe=teal!35!black,title={Before you move on}]
What does soigner mean? (\textquotedblleft{}To look after, to treat, to nurse\textquotedblright{}.) Say it once more.
\end{tcolorbox}
